\documentclass[11pt]{amsart}
\usepackage{amsmath,amssymb,amsthm}
\usepackage[margin=1.1in]{geometry}
\usepackage{hyperref}

\newtheorem{theorem}{Theorem}[section]
\newtheorem{lemma}[theorem]{Lemma}
\newtheorem{proposition}[theorem]{Proposition}
\theoremstyle{remark}
\newtheorem{remark}[theorem]{Remark}

\newcommand{\cA}{\mathcal{A}}
\newcommand{\cB}{\mathcal{B}}
\newcommand{\cE}{\mathcal{E}}
\newcommand{\cP}{\mathcal{P}}
\newcommand{\cT}{\mathcal{T}}
\newcommand{\Sg}{\mathfrak{S}}
\newcommand{\verify}[1]{\textbf{[verify: #1]}}

\title[Sums of two almost prime sums of two squares]{Even integers as sums of two sums of two squares
with few prime factors}
\author{Ruqing Chen}
\address{GUT Geoservice Inc., Montreal, Canada}
\email{ruqing@hotmail.com}

\begin{document}

\begin{abstract}
DRAFT. Let $N\equiv 2\pmod 4$ be sufficiently large. We show that the number of $n\le N$ such that $n$
and $N-n$ are both sums of two squares and $\Omega(n(N-n))\le 12$ is at least
$2\,\Sg(N)N/(\log N)^2$, where $\Sg(N)$ is the binary Goldbach singular series. The level of distribution
of the sequence $n(N-n)$ is $1$, but sifting the primes $q\equiv 3\pmod 4$ up to $N^{1/2}$ with the linear
sieve falls exactly on the sieving limit $s=2$. We therefore sift these primes only up to $N^{1/2-\delta}$
and remove the survivors with two large prime factors $\equiv 3\pmod 4$ on the same side by a switching
argument at the Bombieri--Vinogradov level. The final numerical inequality is verified with interval
arithmetic.
\end{abstract}

\maketitle

%=====================================================================================================
\section{Introduction}

Every positive integer is a sum of four squares, so every $N$ is of the form $n_1+n_2$ with $n_1,n_2$
sums of two squares. We ask for such representations in which $n_1$ and $n_2$ have few prime factors.
This is a two-sided analogue of the problem treated in \cite{Chen1}, where one summand is a prime and
the other an almost prime sum of two squares, and of the classical problem of pairs $n$, $n+h$ of sums of
two squares \cite{Hooley,Indlekofer}.

\begin{theorem}\label{thm:main}
For every sufficiently large $N\equiv 2\pmod 4$,
\begin{multline*}
\#\{n\le N:\ n\equiv 1\ (\mathrm{mod}\ 4),\ n=x^2+y^2,\ N-n=z^2+w^2,\ \Omega(n(N-n))\le 12\}\\
\ge\ 2\,\Sg(N)\,\frac{N}{(\log N)^2},
\end{multline*}
where $\Sg(N)=2\prod_{\ell>2}\bigl(1-(\ell-1)^{-2}\bigr)\prod_{\ell\mid N,\,\ell>2}\frac{\ell-1}{\ell-2}$.
\end{theorem}

\begin{remark}[The prime $2$]
The restriction $N\equiv2\pmod4$ is natural. The odd part of a sum of two squares is $\equiv1\pmod 4$.
If $N\equiv0\pmod 4$, both summands must contain powers of $2$, and for $N=2^j$ the only representation is
$n_1=n_2=2^{j-1}$, so $\Omega(n_1n_2)$ is unbounded. For general $N$ one should count prime factors of the
odd parts only; we do not pursue this here.
\end{remark}

\begin{remark}
Floating point computations (not certified) indicate that the same method gives $\Omega(n(N-n))\le 11$
when $\delta$ is taken close to $1/4$.
\end{remark}

\subsection*{Method}
Restrict to $n\equiv1\pmod4$; then $N-n\equiv1\pmod4$. If neither $n$ nor $N-n$ has a prime factor
$q\equiv3\pmod 4$, both are sums of two squares. The sequence $\{n(N-n)\}$ has level of distribution
$N^{1-\varepsilon}$ (Lemma~\ref{lem:level}), and the primes $q\equiv3\pmod4$ have two roots of
$n(N-n)\equiv0\pmod q$, so they form a sifting set of dimension $1$. Sifting them up to $N^{1/2}$ at level
$N$ gives $s=2$, where the linear lower bound function vanishes. We sift them up to $y=N^{1/2-\delta}$ instead.
Since $n\equiv N-n\equiv1\pmod 4$, each side has an even number of such prime factors, so a survivor that is
not a pair of sums of two squares has two large prime factors $\equiv3\pmod4$ on the \emph{same} side; these
are counted from above by switching (Section~\ref{sec:switch}). Prime factors are controlled by Richert's
weights applied to $n(N-n)$ jointly.

\begin{remark}
The literature has not yet been checked systematically for this statement.
\end{remark}

%=====================================================================================================
\section{The sequence $n(N-n)$}\label{sec:level}

Let $N\equiv 2\pmod 4$ be large, $\cP_1=\{q\equiv1\pmod4\}$, $\cP_3=\{q\equiv3\pmod4\}$, and
\[
\cA=\{\,a_n=n(N-n):\ 1\le n\le N,\ n\equiv1\ (\mathrm{mod}\ 4)\,\},\qquad X=N/4.
\]
For odd squarefree $d$ let $\rho(d)$ be the number of solutions of $n(N-n)\equiv0\pmod d$; it is
multiplicative with $\rho(q)=2$ for $q\nmid N$ and $\rho(q)=1$ for $q\mid N$. Put $g(d)=\rho(d)/d$, so that
\[
g(q)=\frac2q\quad(q\nmid N),\qquad g(q)=\frac1q\quad(q\mid N),
\]
and define $r_d$ by $\#\{n:\ a_n\in\cA,\ d\mid a_n\}=Xg(d)+r_d$.

\begin{lemma}[Level of distribution]\label{lem:level}
$|r_d|\le\rho(d)$, and for $D\le N$,
$\sum_{d\le D}\mu^2(d)\tau_3(d)|r_d|\ll D(\log D)^{C}$. In particular the level of distribution is
$N^{1-\varepsilon}$.
\end{lemma}

\begin{proof}
The condition $d\mid a_n$ together with $n\equiv1\pmod 4$ restricts $n$ to $\rho(d)$ classes modulo $4d$,
each containing $N/(4d)+O(1)$ integers $n\le N$. Hence $|r_d|\le\rho(d)\le 2^{\omega(d)}$, and the sum is
$\ll\sum_{d\le D}\tau_3(d)2^{\omega(d)}\ll D(\log D)^{C}$.
\end{proof}

\begin{lemma}[Local factors]\label{lem:local}
As $z\to\infty$,
\[
\prod_{2<q<z}(1-g(q))\sim\frac{\Sg(N)}{2}\Bigl(\frac{2e^{-\gamma}}{\log z}\Bigr)^2 .
\]
\end{lemma}

\begin{proof}
For $q\nmid N$, $(1-2/q)(1-1/q)^{-2}=1-(q-1)^{-2}$; for $q\mid N$, $(1-1/q)(1-1/q)^{-2}=q/(q-1)$. The product
of these factors over odd $q$ converges to $\Sg(N)/2$, and $\prod_{2<q<z}(1-1/q)\sim 2e^{-\gamma}/\log z$.
\end{proof}

We use the following sieve tools; see \cite{Chen1} for statements and references.
The linear sieve ($\kappa=1$, $\beta=2$) has functions $F_1$, $f_1$ with $F_1(s)=2e^\gamma/s$ for $0<s\le3$ and
$f_1(s)=2e^{\gamma}\log(s-1)/s$ for $2\le s\le 4$, $f_1(s)=0$ for $s\le2$
\cite{IwaniecRosser}; the semi-linear sieve ($\kappa=1/2$, $\beta=1$) has functions $F_{1/2}$,
$f_{1/2}$ \cite{IwaniecHalf}. The vector sieve inequality \cite{BF} is
$ab\ge a^-b^++a^+b^--a^+b^+$ for $a,b\in\{0,1\}$, $a^-\le a\le a^+$, $b^-\le b\le b^+$, $a^+,b^+\ge0$.
The Bombieri--Vinogradov theorem for convolutions is used as in \cite[Lemma~2.4]{Chen1}.

%=====================================================================================================
\section{The weighted sieve and switching}\label{sec:switch}

Fix $0<\alpha<\delta<1/4$, $0<\lambda$, $\alpha<\beta_1\le1/2$, and put
\[
y=N^{1/2-\delta},\quad z=N^{\alpha},\quad y_1=N^{\beta_1},\quad
P_I=\prod_{\substack{q\in\cP_3\\ q<y}}q,\quad P_S=\prod_{\substack{q\in\cP_1\\ q<z}}q,\quad P=P_IP_S .
\]
Every odd prime $q<z$ divides $P$. For $a_n\in\cA$ with $(a_n,P)=1$ define
\[
w_n=1-\lambda\sum_{\substack{q\mid a_n,\ q\in\cP_1\\ z\le q<y_1}}\Bigl(1-\frac{\log q}{\log y_1}\Bigr),
\qquad W=\sum_{(a_n,P)=1}w_n .
\]

\begin{lemma}[Structure]\label{lem:structure}
Let $(a_n,P)=1$. Then either
\begin{enumerate}
\item[(i)] $n$ and $N-n$ are both sums of two squares, all their prime factors are $\ge z$, and those
in $\cP_3$ occur only as squares $q^2$; or
\item[(ii)] for $m'=n$ or $m'=N-n$ we have $m'=q_1q_2m$ with $q_1<q_2$, $q_1,q_2\in\cP_3$, $q_1\ge y$, and
$m\le N^{2\delta}$ composed of primes of $\cP_1$ that are $\ge z$.
\end{enumerate}
\end{lemma}

\begin{proof}
Each of $n$, $N-n$ is odd and $\equiv1\pmod4$, so each has an even number of prime factors in $\cP_3$,
counted with multiplicity. They are all $\ge y$, and $y^4>N$, so each side has $0$ or $2$ of them. If on both
sides there are none, or they coincide, both sides are sums of two squares. Otherwise some side is as in
(ii). In particular, two exceptional prime factors never lie on different sides.
\end{proof}

\begin{lemma}[Prime factors]\label{lem:omega}
If $a_n$ is as in Lemma~\ref{lem:structure}(i), squarefree, and $w_n>0$, then
$\Omega(n(N-n))<1/\lambda+2/\beta_1$.
\end{lemma}

\begin{proof}
All prime factors of $a_n$ are in $\cP_1$ and $\ge z$, and $\log a_n\le 2\log N$. As in
\cite[Lemma~3.2]{Chen1},
$\sum_{q\mid a_n,\,q<y_1}(1-\log q/\log y_1)\ge\Omega(a_n)-\log a_n/\log y_1\ge\Omega(a_n)-2/\beta_1$.
\end{proof}

Non-squarefree $a_n$ with a square factor $q^2$, $q\ge z$, number $\ll N^{1-\alpha}$ and are negligible.

\begin{proposition}[Main term]\label{prop:W}
Let $\theta_I\in(0,1)$ and, for $u\in[\alpha,\beta_1]$, $t(u)\in(0,1-u)$. Put $s_I=\theta_I/(1/2-\delta)$,
$s_S=(1-\theta_I)/\alpha$,
\[
L=f_1(s_I)F_1(s_S)+F_1(s_I)f_1(s_S)-F_1(s_I)F_1(s_S),
\]
\[
R=\int_\alpha^{\beta_1}\Bigl(1-\frac u{\beta_1}\Bigr)F_1\Bigl(\frac{t(u)}{1/2-\delta}\Bigr)
F_1\Bigl(\frac{1-u-t(u)}{\alpha}\Bigr)\frac{du}{u}.
\]
Then $W\ge(L-\lambda R+o(1))\,XV$, where $V=\prod_{q\mid P}(1-g(q))$ and
\[
XV\sim\frac{e^{-2\gamma}}{2\alpha(1/2-\delta)}\;\Sg(N)\frac{N}{(\log N)^2}.
\]
\end{proposition}

\begin{proof}
Apply the vector sieve with $a=1_{(a_n,P_I)=1}$ and $b=1_{(a_n,P_S)=1}$. Both are linear sieves:
$P_I$ and $P_S$ consist of primes from one half of the residue classes modulo $4$, each with
$\rho(q)=2$ (apart from the finitely many $q\mid N$ below a given bound, which do not affect the dimension).
Use levels $N^{\theta_I}$ and $N^{1-\theta_I-\varepsilon}$; the combined moduli are $\le N^{1-\varepsilon}$
and Lemma~\ref{lem:level} controls the remainders. For the weights, $q\mid a_n$ for $q\in\cP_1$, $q\nmid N$,
has density $g(q)=2/q$, and $\sum_{q\in\cP_1}2/q$ over $N^u\le q<N^{u+du}$ is $du/u$. The upper bound vector
sieve for $\cA_q$ at levels $N^{t(u)}$, $N^{1-u-t(u)-\varepsilon}$ gives $R$. Finally
$V=\prod_{2<q<z}(1-g(q))\prod_{q\in\cP_3,\,z\le q<y}(1-g(q))$, and Lemma~\ref{lem:local} together with
$\prod_{q\in\cP_3,\,z\le q<y}(1-2/q)\sim\log z/\log y$ gives the asymptotic for $XV$.
\end{proof}

\begin{proposition}[Switching]\label{prop:switch}
Let $\cE$ be the set of $n$ with $(a_n,P)=1$ in case (ii) of Lemma~\ref{lem:structure}. Then
\[
\#\cE\le\Bigl(4\,c(\delta,\alpha)\,e^{\gamma}\sqrt{\alpha(1/2-\delta)}\;U+o(1)\Bigr)\,XV,
\]
where $c(\delta,\alpha)$ is the switching density of \cite[Prop.~4.1]{Chen1} and
\[
U=\min_{0<t<1/2}F_{1/2}\Bigl(\frac{t}{1/2-\delta}\Bigr)F_{1/2}\Bigl(\frac{1/2-t}{\alpha}\Bigr).
\]
\end{proposition}

\begin{proof}
Write $\cE=\cE_1\cup\cE_2$ according as the exceptional pair lies in $n$ or in $N-n$; the two cases are
symmetric, which accounts for a factor $2$. Consider $\cE_1$. Let $\cT$ be the set of triples
$(q_1,q_2,m)$ as in Lemma~\ref{lem:structure}(ii) with $t=q_1q_2m\le N$, so that
$|\cT|=(c(\delta,\alpha)+o(1))N/\log N$ as in \cite{Chen1}; note $t\equiv1\pmod4$ automatically.
Every $n\in\cE_1$ equals such a $t$, and $N-t$ has no prime factor in $P$. Hence $\#\cE_1$ is at most the
number of elements of $\cB=\{N-t:\ t\in\cT\}$ coprime to $P$.

Let $\ell\mid P$. No prime factor of $t$ divides $P$ (the $q_i$ are $\ge y$ and the prime factors of $m$
are split and $\ge z$), so $t$ is equidistributed in the nonzero classes modulo $\ell$. Hence, for
$\ell\nmid N$, $\ell\mid N-t$ with density $g_\cB(\ell)=1/(\ell-1)$, while for $\ell\mid N$ we have
$\ell\nmid N-t$, so $g_\cB(\ell)=0$. The upper bound vector sieve for $\cB$ with the semi-linear components
$P_I$ and $P_S$ at total level $N^{1/2}(\log N)^{-B}$ (Bombieri--Vinogradov for the convolution
$q_1*(q_2m)$) gives
\[
\#\cE_1\le\bigl(U+o(1)\bigr)\,|\cB|\prod_{q\mid P}(1-g_\cB(q)).
\]
Now $(1-1/(\ell-1))(1-1/\ell)^{-1}=1-(\ell-1)^{-2}$ and $(1-0)(1-1/\ell)^{-1}=\ell/(\ell-1)$ for $\ell\mid N$:
these are the same local factors as in Lemma~\ref{lem:local}. Hence
\[
\prod_{q\mid P}(1-g_\cB(q))\sim\frac{\Sg(N)}{2}\cdot\frac{2e^{-\gamma}}{\log z}\Bigl(\frac{\log z}{\log y}\Bigr)^{1/2},
\]
and therefore
\[
\#\cE_1\le\bigl(U+o(1)\bigr)\,c\,\frac{N}{\log N}\cdot\frac{\Sg(N)e^{-\gamma}}{\alpha\log N}
\Bigl(\frac{\alpha}{1/2-\delta}\Bigr)^{1/2}
=\bigl(2c\,e^{\gamma}\sqrt{\alpha(1/2-\delta)}\,U+o(1)\bigr)XV,
\]
using Proposition~\ref{prop:W} for $XV$. The singular series cancels exactly. Adding $\cE_2$ gives the claim.
\end{proof}

%=====================================================================================================
\section{Numerical verification and proof of Theorem~\ref{thm:main}}\label{sec:numerics}

By Lemma~\ref{lem:structure} and $w_n\le1$, the weighted count over case (i) is at least
$W-\#\cE\ge(\mathrm{net}+o(1))XV$ with
\[
\mathrm{net}=L-\lambda R-4c(\delta,\alpha)e^{\gamma}\sqrt{\alpha(1/2-\delta)}\,U .
\]
We take $\delta=0.22$, $\alpha=0.01$, $\lambda=1/8$, $\beta_1=2/5$, so that $1/\lambda+2/\beta_1=13$ and,
by Lemma~\ref{lem:omega}, $\Omega(n(N-n))\le12$ for the counted $n$.

\begin{lemma}\label{lem:numerics}
For these parameters (with $\theta_I=0.9495$), $L\ge0.913373$, $R\le3.144708$, $c(\delta,\alpha)\le1.092388$,
$U\le1.166741$ and $\mathrm{net}\ge0.039808$.
\end{lemma}

\begin{proof}[Method of verification]
The script \texttt{code/certify.py} uses outward-rounded interval arithmetic as in \cite[Section~5]{Chen1},
now also for the linear sieve functions, whose delay-differential equations are integrated on the grid
$10^{-3}\mathbb Z$ up to $s=110$ with monotone enclosures.
\end{proof}

\begin{proof}[Proof of Theorem~\ref{thm:main}]
By Lemma~\ref{lem:numerics} and Proposition~\ref{prop:W}, the number of counted $n$ is at least
$(0.0398+o(1))\frac{e^{-2\gamma}}{2\cdot0.01\cdot0.28}\Sg(N)\frac{N}{(\log N)^2}
\ge 2\,\Sg(N)\frac{N}{(\log N)^2}$ for large $N$ (the numerical constant is $\ge2.24$).
\end{proof}

\begin{thebibliography}{9}
\bibitem{BF} J.~Br\"udern, \'E.~Fouvry, \emph{Lagrange's four squares theorem with almost prime variables},
J. reine angew. Math. \textbf{454} (1994), 59--96.
\bibitem{Chen1} R.~Chen, \emph{On primes $p$ for which $N-p$ is an almost prime sum of two squares, and other
class number one forms}, preprint, doi:10.5281/zenodo.23136112.
\bibitem{Hooley} C.~Hooley, \emph{On the intervals between numbers that are sums of two squares}
\verify{exact reference}.
\bibitem{Indlekofer} K.-H.~Indlekofer, \emph{Scharfe untere Absch\"atzung f\"ur die Anzahlfunktion der
B-Zwillinge}, Acta Arith. \textbf{26} (1974/75), 207--212 \verify{check}.
\bibitem{IwaniecHalf} H.~Iwaniec, \emph{The half dimensional sieve}, Acta Arith. \textbf{29} (1976), 69--95.
\bibitem{IwaniecRosser} H.~Iwaniec, \emph{Rosser's sieve}, Acta Arith. \textbf{36} (1980), 171--202.
\end{thebibliography}

\end{document}
